

\section{Introduction} \label{sec:introduction}
% X: What are we trying to do and why is it relevant?


\begin{figure}[t]
    \centering
    \begin{subfigure}[b]{0.485\linewidth}
        \centering
        \includegraphics[width=\linewidth]{figures/eval_aware_digram.pdf}
        \caption{Obfuscated model evades monitor}
        \label{fig:sub1}
    \end{subfigure}
    \hfill
    \begin{subfigure}[b]{0.485\linewidth}
        \centering
        \includegraphics[width=\linewidth]{figures/std_model.pdf}
        \caption{Standard backdoored model detected}
        \label{fig:sub2}
    \end{subfigure}
    \caption{\textbf{Evaluation-aware models can evade latent-space monitors.} (a) An obfuscated model shifts harmful representations outside the monitor's detection scope through feature migration, successfully evading detection. (b) A standard backdoored model's harmful representations remain within the monitor's scope and are detected. Our work characterizes these evasion strategies and introduces \textsc{SafetyNet}, an ensemble detector that monitors multiple geometric subspaces to prevent such evasion.}
    \label{fig:main}
\end{figure}


% X: What are we trying to do and why is it relevant?
As language models become more capable, their evaluation awareness increases \cite{chaudhary2025evaluation, nguyen2025probingsteeringevaluationawareness, apollo2025claude, needham2025largelanguagemodelsknow}---behaving differently when monitored or evaluated.
This gives them the capability to detect whether they are being evaluated or monitored and hence hide certain goals.
This is a critical problem if these goals are harmful---embedded unintentionally~\cite{gillen2024misalignmentcatastrophe} or intentionally~\cite{hubinger2024sleeper} during pre-training or fine-tuning---as the models will be able to conceal them during monitoring.
If successful in fooling monitors, these models can pursue such goals during deployment, leveraging their ability to accomplish long tasks~\cite{metr_long_tasks} in an agentic fashion---enabling scheming~\cite{incontext_Scheming}, sandbagging~\cite{vanderweij2025aisandbagginglanguagemodels, benton2024sabotageevaluationsfrontiermodels}, and pursuit of misaligned goals with catastrophic consequences~\cite{gillen2024misalignmentcatastrophe}.


% Y: Why is this hard?
Therefore, building effective safety monitors faces fundamental challenges. 
First, pre-deployment testing cannot anticipate all harmful behaviors~\citep{metr-s-gpt-4-5-pre-deployment-evaluations, apolloresearchsafetytest, anthropicsafetytest}, especially as models become more capable.
Second, evaluation-aware models can strategically modify their behavior during testing~\citep{berglund2023taken,laine2024myself,chaudhary2025evaluation}, appearing safe while concealing harmful capabilities.
Third, even latent-space monitors are vulnerable: models could learn to evade detection by shifting representations between geometric spaces that individual detectors cannot capture.
This creates an adversarial arms race where increasingly capable models develop more effective evasion strategies, making it essential to first understand \textit{how} models can fool latent-space monitors before building monitors that remain robust against such evasion.

% Z: How do we solve it (our contribution!)
We address this through a two-part approach.
First, we characterize how models evade latent-space monitors through red-team experiments that synthetically train models to fool detection systems using obfuscation objectives.
This reveals two key evasion mechanisms: \textit{geometric shifting} between linear and non-linear subspaces (under similarity loss obfuscation, autoencoder accuracy falls to 69\% on Llama-3 while PCA rises to 98\%), and \textit{covariate manipulation} that alters feature relationships.
Second, leveraging these insights, we introduce \textsc{SafetyNet}, a robust ensemble detector that monitors multiple geometric properties---PCA (linear), AE/VAE (non-linear manifolds), and Mahalanobis distance (covariance)---preventing models from hiding in any single representational space (Figure~\ref{fig:main} and Figure~\ref{fig:methodology}).
We provide theoretical analysis showing that evading all detectors simultaneously requires either degrading task performance or leaving detectable signatures in at least one geometric subspace (Section~\ref{sec:method}).

Importantly, while our red-team approach modifies model weights through fine-tuning, the evasion strategies we identify remain relevant for real-world deployment.
Modern language models can modify their representational behavior through in-context learning~\cite{dai2023gptlearnincontextlanguage} and chain-of-thought reasoning~\cite{huang2025transformerslearnimplementmultistep}, which effectively alter the functional composition of representations during inference without explicit weight updates.
These inference-time adaptations can achieve similar geometric shifts to those we study through weight modifications, making our findings applicable even when models cannot directly access or modify their parameters.


% 1: How do we verify that we solved it (experiments and results)
Across five model families, including Llama-2 7B, Llama-3 8B, Qwen-2.5 3B, Gemma 7B, and Mistral 7B, on MAD \cite{MAD} and Sleeper Agent benchmarks~\cite{hubinger2024sleeper}, \textsc{SafetyNet} achieves 98.9\% average detection accuracy versus 3.6\% for CROW~\cite{crow} and 0.4\% for Beatrix~\cite{maBeatrixResurrectionsRobust2022}.
Against adaptive obfuscation attacks designed to evade detection, \textsc{SafetyNet} maintains 89.8\% accuracy while baseline methods collapse below 10\%.
Causal intervention experiments validate that our detectors capture causally relevant features: manipulating high-activation attention vectors produces significant logit differences, confirming detection signals drive model behavior rather than reflecting spurious correlations (Section~\ref{sec:causal_validation}).

\begin{figure}[t]
    \centering
    \input{table/new_methodology_fig}
    \caption{\textbf{\textsc{SafetyNet} ensemble architecture and adversarial training framework.}
    \textsc{SafetyNet} combines four complementary \textcolor{ForestGreen!60}{detectors}---Variational Autoencoder (VAE), Autoencoder (AE), PCA, and Mahalanobis Distance---to classify model internal states as \textcolor{NavyBlue}{benign} or \textcolor{violetcolor}{harmful} with probability $P(Y)$.
    Each detector monitors different geometric properties: VAE and AE capture non-linear manifold structure, PCA captures linear subspaces, and Mahalanobis distance captures covariance relationships.
    The framework employs outlier detection, treating harmful representations as out-of-distribution relative to benign samples.
    In our \textcolor{violetcolor}{red-team}/\textcolor{NavyBlue}{blue-team} evaluation, the red team trains models to evade detection by manipulating attention layer representations, while the blue team builds robust ensemble monitors that remain effective against such adaptive attacks.}
    \label{fig:methodology}
\end{figure}